\documentclass[10pt,a4paper]{amsart}
\usepackage[margin=19mm]{geometry}
\usepackage{amsmath,amssymb,mathtools}
\usepackage[scr=rsfs,cal=euler]{mathalfa}
\usepackage{xfrac}
\usepackage[unicode,hidelinks,bookmarks=false]{hyperref}
\newcommand{\ie}{i.e.\ }
\newcommand{\cf}{cf.\ }
\newcommand{\resp}{respectively\ }
\newcommand{\Subsection}{\subsection}
\providecommand{\nobreakdash}{\nobreak\hbox{-}}
\setlength{\emergencystretch}{2em}
\multlinegap0pt
\usepackage{bm}
\usepackage{bbm}
\usepackage{dsfont}
\usepackage{xspace}
\usepackage{cancel}
\usepackage{mathtools}
\usepackage{amsthm}
\makeatletter
\@ifundefined{defn}{\newtheorem{defn}{Defn}}{}
\@ifundefined{lem}{\newtheorem{lem}[defn]{Lemma}}{}
\makeatother
\newcommand\Arg{\operatorname{Arg}}
\newcommand\Ima{\operatorname{Im}\,}
\newcommand\PP{\mathbb P}
\newcommand\Rea{\operatorname{Re}\,}
\newcommand\cA{\mathcal A}
\newcommand\cO{\mathcal O}
\newcommand\fD{\mathfrak D}
\newcommand\iin{\mathrm{in}}
\title[Scattering diagrams, stability conditions, and coherent shea]{Scattering diagrams, stability conditions, and coherent sheaves on $\mathbb{P}^2$}
\author{Pierrick Bousseau}
\date{Source: arXiv:1909.02985v3 (CC BY 4.0)}
\begin{document}
\maketitle

\noindent\emph{Source and licence.} Scattering diagrams, stability conditions, and coherent sheaves on $\mathbb{P}^2$. Version arXiv:1909.02985v3 (CC BY 4.0), distributed under the Creative Commons Attribution 4.0 International licence, as recorded in the arXiv licence metadata for this exact version and shown on the abstract page https://arxiv.org/abs/1909.02985v3. Full rights notice: licenses/arxiv-1909.02985-CC-BY-4.0.txt.\par\smallskip

\noindent\textit{Editorial scope.} Counted unit: the Lem whose source label is \texttt{lem\_coincide\_int\_T} in the article (source0.tex lines 4540--4547, source label \texttt{lem\_coincide\_int\_T}), delivered together with the article's own complete original proof of it (source0.tex lines 4549--4863, one \texttt{proof} environment of 315 lines). No mathematics is written, repaired or re-derived: statement and proof are the article's own text line for line. Required context retained uncounted: none. Every definition, display and label of those lines is retained unchanged. Omitted from this package and not redistributed: 1--4539, 4548--4548, 4864--7175. Nothing inside a retained excerpt is omitted or altered: every retained body is delivered exactly as the article's lines are written, so no omission receipt of any kind accompanies this package. Citations: the retained excerpts carry no citation command at all, so no bibliography entry of the article is transcribed and no reference is redistributed. Reference adaptation: none was needed and none was made; every reference inside the delivered unit resolves inside it, so no unresolved reference or citation is produced. Printed slips kept literal: no printed word, symbol, hypothesis or number of the counted statement or of its proof was altered. Layout and class: the article's own class is \texttt{amsart} with packages including babel, inputenc, amsmath, graphicx, todonotes, amssymb, amscd, latexsym, epsfig, xy, euscript, amsthm, tikz, MnSymbol; the wrapper is the lane's pinned \texttt{amsart} editorial wrapper at 10pt on A4, which restates the theorem-like environments the retained text needs and adds the article's own macro definitions that the retained text needs (\texttt{\textbackslash Arg}, \texttt{\textbackslash Ima}, \texttt{\textbackslash PP}, \texttt{\textbackslash Rea}, \texttt{\textbackslash cA}, \texttt{\textbackslash cO}, \texttt{\textbackslash fD}, \texttt{\textbackslash iin}), with extra wrapper packages (bm, bbm, dsfont, xspace, cancel, mathtools). The delivered numbering is the unit's own. Assets: no retained span contains a figure environment, a graphics-inclusion command, a PDF-page inclusion, a table or a TikZ picture; no image file is copied into this package, the package declares no asset dependency and no image file is redistributed with it. Licence: version arXiv:1909.02985v3 of this article is distributed under the Creative Commons Attribution 4.0 International licence, as recorded in the arXiv licence metadata for this exact version and shown on the abstract page https://arxiv.org/abs/1909.02985v3. A full rights notice is delivered at licenses/arxiv-1909.02985-CC-BY-4.0.txt. Characters: every retained excerpt is pure ASCII, so the delivered engine, XeLaTeX with its default Latin Modern text font, reports no missing character.
\begin{lem} \label{lem_coincide_int_T}
The scattering diagram $\fD^{\PP^2}_{u,v}$ is empty in restriction to the interior of 
$\bar{U}^\iin_{0,T}$.
In other words, the scattering diagrams 
$\fD^{\PP^2}_{u,v}$ and 
$S(\fD^\iin_{u,v})$ coincide in restriction to the interior of 
$\bar{U}^\iin_{0,T}$.
\end{lem}

\begin{proof}
Let $\sigma$ be a point in the interior of $\bar{U}^\iin_{0,T}$.
Then $\cO(-1)[1]$ and $\cO(1)$ belong to 
$\cA^\sigma$, with 
\[\Rea Z^\sigma_{\gamma(\cO(-1)[1])}=-\Rea Z^\sigma_{\gamma(\cO(-1))}>0\,,\]
and
\[\Rea Z^\sigma_{\gamma(\cO(1))}<0\,.\]
We also have 
\[ \Rea Z^\sigma_{\gamma(\cO)}>0\,,\]
and, depending if $x \geqslant 0$ or 
$x<0$, $\cO$ belongs to $\cA^\sigma$ or 
$\cO[1]$ belongs to 
$\cA^\sigma$.
In any case, the objects $\cO(-1)[2]$, $\cO[1]$, and $\cO(1)$ belong to $\cA^\sigma[Z^\sigma,1/2]$, and so 
$\cA^\sigma[Z^\sigma, 1/2]=\cA_0$.
On the other hand, we have 
\[ \Rea Z_{\gamma(\cO(-1)[2])}^\sigma 
<0\,, \Rea Z_{\gamma(\cO[1])}^\sigma<0 \,,
\Rea Z_{\gamma(\cO(1))}^\sigma <0\,.\]





As $\cA_0$ is a category of quiver representations, with simple objects 
$\cO(-1)[2]$, $\cO[1]$ and $\cO(1)$, we have that, for every $E$ nonzero object of 
$\cA^\sigma[Z^\sigma,1/2]$, the central charge
$Z^\sigma[1/2](E)$ is contained in the cone of linear combinations with 
nonnegative coefficients of 
$Z^\sigma[1/2]_{\gamma(\cO(-1))[2])}$,
$Z^\sigma[1/2]_{\gamma(\cO[1])}
$, 
$Z^\sigma[1/2]_{\gamma(\cO(1))}$.
As $Z^\sigma[1/2]_\gamma=-i Z^\sigma_\gamma$ for every $\gamma \in \Gamma$, 
we have  
\[ \Ima Z^\sigma[1/2]_{\gamma(\cO(-1))[2])}>0\,,
\Ima Z^\sigma[1/2]_{\gamma(\cO[1])}>0
\,,
\Ima Z^\sigma[1/2]_{\gamma(\cO(1))}>0\,.\]
and so
we conclude that 
\[\Ima Z^\sigma[1/2](E)>0\,.\]
Thus, if $E$ is a $\sigma$-semistable object of $\cA^\sigma$, then 
\begin{itemize}
\item[(i)] either, $\frac{1}{\pi} \Arg Z^\sigma(E)>\frac{1}{2}$, and so 
$\Rea Z^\sigma(E)<0$,
\item[(ii)] or, $\frac{1}{\pi} \Arg Z^{\sigma}(E) 
\leqslant \frac{1}{2}$, and so $E[1]$
is a $\sigma[1/2]$-semistable object 
of $\cA^\sigma[Z^\sigma, \frac{1}{2}]$, so
$\Ima Z^\sigma[1/2](E[1])>0$, and so
$\Rea Z^\sigma(E)>0$.
\end{itemize}
In any case, we have $\Rea Z^\sigma(E) \neq 0$. It follows that $\fD^{\PP^2}_{u,v}$, restricted to the interior of 
$\bar{U}_{0,T}^{\iin}$, is empty, as 
$S(\fD^\iin_{u,v})$.

Figure: an example of configuration of central charges for $\sigma=(x,y)$ in the interior of $\bar{U}^\iin_{0,T}$ with $x <0$. If $E$ is a $\sigma$-semistable object of $\cA^\sigma$, then $Z^\sigma(E)$ belongs to the dotted region.
\begin{center}
\setlength{\unitlength}{1.2cm}
\begin{picture}(10,6)
\put(5,3){\line(1,0){4}}
\put(5,3){\line(-1,0){4}}
\put(5,3){\line(0,1){3}}
\put(5,3){\line(0,-1){3}}
\put(5,3){\vector(1,1){2}}
\put(7.1,5){$Z^\sigma_{\gamma(\cO(-1)[1])}$}
\put(5,3){\vector(-1,-1){2}}
\put(2.5,0.8){$Z^\sigma_{\gamma(\cO(-1)[2])}$}
\put(5,3){\vector(-1,3){1}}
\put(3,5.5){$Z^\sigma_{\gamma(\cO(1))}$}
\put(5,3){\vector(-2,1){2}}
\put(1.8,4){$Z^\sigma_{\gamma(\cO[1])}$}
\put(5.25,3.25){\circle*{0.05}}
\put(5.5,3.25){\circle*{0.05}}
\put(5.75,3.25){\circle*{0.05}}
\put(6,3.25){\circle*{0.05}}
\put(6.25,3.25){\circle*{0.05}}
\put(6.5,3.25){\circle*{0.05}}
\put(6.75,3.25){\circle*{0.05}}
\put(7,3.25){\circle*{0.05}}
\put(7.25,3.25){\circle*{0.05}}
\put(7.5,3.25){\circle*{0.05}}
\put(7.75,3.25){\circle*{0.05}}
\put(8,3.25){\circle*{0.05}}
\put(8.25,3.25){\circle*{0.05}}
\put(8.5,3.25){\circle*{0.05}}
\put(5.5,3.5){\circle*{0.05}}
\put(5.75,3.5){\circle*{0.05}}
\put(6,3.5){\circle*{0.05}}
\put(6.25,3.5){\circle*{0.05}}
\put(6.5,3.5){\circle*{0.05}}
\put(6.75,3.5){\circle*{0.05}}
\put(7,3.5){\circle*{0.05}}
\put(7.25,3.5){\circle*{0.05}}
\put(7.5,3.5){\circle*{0.05}}
\put(7.75,3.5){\circle*{0.05}}
\put(8,3.5){\circle*{0.05}}
\put(8.25,3.5){\circle*{0.05}}
\put(8.5,3.5){\circle*{0.05}}
\put(5.75,3.75){\circle*{0.05}}
\put(6,3.75){\circle*{0.05}}
\put(6.25,3.75){\circle*{0.05}}
\put(6.5,3.75){\circle*{0.05}}
\put(6.75,3.75){\circle*{0.05}}
\put(7,3.75){\circle*{0.05}}
\put(7.25,3.75){\circle*{0.05}}
\put(7.5,3.75){\circle*{0.05}}
\put(7.75,3.75){\circle*{0.05}}
\put(8,3.75){\circle*{0.05}}
\put(8.25,3.75){\circle*{0.05}}
\put(8.5,3.75){\circle*{0.05}}
\put(6,4){\circle*{0.05}}
\put(6.25,4){\circle*{0.05}}
\put(6.5,4){\circle*{0.05}}
\put(6.75,4){\circle*{0.05}}
\put(7,4){\circle*{0.05}}
\put(7.25,4){\circle*{0.05}}
\put(7.5,4){\circle*{0.05}}
\put(7.75,4){\circle*{0.05}}
\put(8,4){\circle*{0.05}}
\put(8.25,4){\circle*{0.05}}
\put(8.5,4){\circle*{0.05}}
\put(6.25,4.25){\circle*{0.05}}
\put(6.5,4.25){\circle*{0.05}}
\put(6.75,4.25){\circle*{0.05}}
\put(7,4.25){\circle*{0.05}}
\put(7.25,4.25){\circle*{0.05}}
\put(7.5,4.25){\circle*{0.05}}
\put(7.75,4.25){\circle*{0.05}}
\put(8,4.25){\circle*{0.05}}
\put(8.25,4.25){\circle*{0.05}}
\put(8.5,4.25){\circle*{0.05}}
\put(6.5,4.5){\circle*{0.05}}
\put(6.75,4.5){\circle*{0.05}}
\put(7,4.5){\circle*{0.05}}
\put(7.25,4.5){\circle*{0.05}}
\put(7.5,4.5){\circle*{0.05}}
\put(7.75,4.5){\circle*{0.05}}
\put(8,4.5){\circle*{0.05}}
\put(8.25,4.5){\circle*{0.05}}
\put(8.5,4.5){\circle*{0.05}}
\put(1.5,3){\circle*{0.05}}
\put(1.5,3.25){\circle*{0.05}}
\put(1.5,3.5){\circle*{0.05}}
\put(1.5,3.75){\circle*{0.05}}
\put(1.5,4){\circle*{0.05}}
\put(1.5,4.25){\circle*{0.05}}
\put(1.5,4.5){\circle*{0.05}}
\put(1.5,4.75){\circle*{0.05}}
\put(1.5,5){\circle*{0.05}}
\put(1.5,5.25){\circle*{0.05}}
\put(1.5,5.5){\circle*{0.05}}
\put(1.5,5.75){\circle*{0.05}}
\put(1.5,6){\circle*{0.05}}
\put(1.75,3){\circle*{0.05}}
\put(1.75,3.25){\circle*{0.05}}
\put(1.75,3.5){\circle*{0.05}}
\put(1.75,3.75){\circle*{0.05}}
%\put(1.75,4){\circle*{0.05}}
%\put(1.75,4.25){\circle*{0.05}}
\put(1.75,4.5){\circle*{0.05}}
\put(1.75,4.75){\circle*{0.05}}
\put(1.75,5){\circle*{0.05}}
\put(1.75,5.25){\circle*{0.05}}
\put(1.75,5.5){\circle*{0.05}}
\put(1.75,5.75){\circle*{0.05}}
\put(1.75,6){\circle*{0.05}}
\put(2,3){\circle*{0.05}}
\put(2,3.25){\circle*{0.05}}
\put(2,3.5){\circle*{0.05}}
\put(2,3.75){\circle*{0.05}}
%\put(2,4){\circle*{0.05}}
%\put(2,4.25){\circle*{0.05}}
\put(2,4.5){\circle*{0.05}}
\put(2,4.75){\circle*{0.05}}
\put(2,5){\circle*{0.05}}
\put(2,5.25){\circle*{0.05}}
\put(2,5.5){\circle*{0.05}}
\put(2,5.75){\circle*{0.05}}
\put(2,6){\circle*{0.05}}
\put(2.25,3){\circle*{0.05}}
\put(2.25,3.25){\circle*{0.05}}
\put(2.25,3.5){\circle*{0.05}}
\put(2.25,3.75){\circle*{0.05}}
%\put(1.75,4){\circle*{0.05}}
%\put(1.75,4.25){\circle*{0.05}}
\put(2.25,4.5){\circle*{0.05}}
\put(2.25,4.75){\circle*{0.05}}
\put(2.25,5){\circle*{0.05}}
\put(2.25,5.25){\circle*{0.05}}
\put(2.25,5.5){\circle*{0.05}}
\put(2.25,5.75){\circle*{0.05}}
\put(2.25,6){\circle*{0.05}}
\put(2.5,3){\circle*{0.05}}
\put(2.5,3.25){\circle*{0.05}}
\put(2.5,3.5){\circle*{0.05}}
\put(2.5,3.75){\circle*{0.05}}
%\put(1.5,4){\circle*{0.05}}
%\put(1.5,4.25){\circle*{0.05}}
\put(2.5,4.5){\circle*{0.05}}
\put(2.5,4.75){\circle*{0.05}}
\put(2.5,5){\circle*{0.05}}
\put(2.5,5.25){\circle*{0.05}}
\put(2.5,5.5){\circle*{0.05}}
\put(2.5,5.75){\circle*{0.05}}
\put(2.5,6){\circle*{0.05}}
\put(2.75,3){\circle*{0.05}}
\put(2.75,3.25){\circle*{0.05}}
\put(2.75,3.5){\circle*{0.05}}
\put(2.75,3.75){\circle*{0.05}}
%\put(1.75,4){\circle*{0.05}}
%\put(1.75,4.25){\circle*{0.05}}
\put(2.75,4.5){\circle*{0.05}}
\put(2.75,4.75){\circle*{0.05}}
\put(2.75,5){\circle*{0.05}}
\put(2.75,5.25){\circle*{0.05}}
\put(2.75,5.5){\circle*{0.05}}
\put(2.75,5.75){\circle*{0.05}}
\put(2.75,6){\circle*{0.05}}
\put(3,3){\circle*{0.05}}
\put(3,3.25){\circle*{0.05}}
\put(3,3.5){\circle*{0.05}}
\put(3,3.75){\circle*{0.05}}
\put(3,4){\circle*{0.05}}
\put(3,4.25){\circle*{0.05}}
\put(3,4.5){\circle*{0.05}}
\put(3,4.75){\circle*{0.05}}
\put(3,5){\circle*{0.05}}
\put(3,5.25){\circle*{0.05}}
%\put(3,5.5){\circle*{0.05}}
%\put(3,5.75){\circle*{0.05}}
\put(3,6){\circle*{0.05}}
\put(3.25,3){\circle*{0.05}}
\put(3.25,3.25){\circle*{0.05}}
\put(3.25,3.5){\circle*{0.05}}
\put(3.25,3.75){\circle*{0.05}}
\put(3.25,4){\circle*{0.05}}
\put(3.25,4.25){\circle*{0.05}}
\put(3.25,4.5){\circle*{0.05}}
\put(3.25,4.75){\circle*{0.05}}
\put(3.25,5){\circle*{0.05}}
\put(3.25,5.25){\circle*{0.05}}
%\put(3.25,5.5){\circle*{0.05}}
%\put(3.25,5.75){\circle*{0.05}}
\put(3.25,6){\circle*{0.05}}
\put(3.5,3){\circle*{0.05}}
\put(3.5,3.25){\circle*{0.05}}
\put(3.5,3.5){\circle*{0.05}}
\put(3.5,3.75){\circle*{0.05}}
\put(3.5,4){\circle*{0.05}}
\put(3.5,4.25){\circle*{0.05}}
\put(3.5,4.5){\circle*{0.05}}
\put(3.5,4.75){\circle*{0.05}}
\put(3.5,5){\circle*{0.05}}
\put(3.5,5.25){\circle*{0.05}}
%\put(3.5,5.5){\circle*{0.05}}
%\put(3.5,5.75){\circle*{0.05}}
\put(3.5,6){\circle*{0.05}}
\put(3.75,3){\circle*{0.05}}
\put(3.75,3.25){\circle*{0.05}}
\put(3.75,3.5){\circle*{0.05}}
\put(3.75,3.75){\circle*{0.05}}
\put(3.75,4){\circle*{0.05}}
\put(3.75,4.25){\circle*{0.05}}
\put(3.75,4.5){\circle*{0.05}}
\put(3.75,4.75){\circle*{0.05}}
\put(3.75,5){\circle*{0.05}}
\put(3.75,5.25){\circle*{0.05}}
%\put(3.75,5.5){\circle*{0.05}}
\put(3.75,5.75){\circle*{0.05}}
\put(3.75,6){\circle*{0.05}}
\put(4,3){\circle*{0.05}}
\put(4,3.25){\circle*{0.05}}
\put(4,3.5){\circle*{0.05}}
\put(4,3.75){\circle*{0.05}}
\put(4,4){\circle*{0.05}}
\put(4,4.25){\circle*{0.05}}
\put(4,4.5){\circle*{0.05}}
\put(4,4.75){\circle*{0.05}}
\put(4,5){\circle*{0.05}}
\put(4,5.25){\circle*{0.05}}
%\put(3,5.5){\circle*{0.05}}
%\put(3,5.75){\circle*{0.05}}
\put(4,6){\circle*{0.05}}
\put(4.25,3){\circle*{0.05}}
\put(4.25,3.25){\circle*{0.05}}
\put(4.25,3.5){\circle*{0.05}}
\put(4.25,3.75){\circle*{0.05}}
\put(4.25,4){\circle*{0.05}}
\put(4.25,4.25){\circle*{0.05}}
\put(4.25,4.5){\circle*{0.05}}
\put(4.25,4.75){\circle*{0.05}}
\put(4.25,5){\circle*{0.05}}
\put(4.25,5.25){\circle*{0.05}}
%\put(4.25,5.5){\circle*{0.05}}
%\put(4.25,5.75){\circle*{0.05}}
%\put(4.25,6){\circle*{0.05}}
\put(4.5,3){\circle*{0.05}}
\put(4.5,3.25){\circle*{0.05}}
\put(4.5,3.5){\circle*{0.05}}
\put(4.5,3.75){\circle*{0.05}}
\put(4.5,4){\circle*{0.05}}
\put(4.5,4.25){\circle*{0.05}}
\put(4.5,4.5){\circle*{0.05}}
\put(4.75,3){\circle*{0.05}}
\put(4.75,3.25){\circle*{0.05}}
\put(4.75,3.5){\circle*{0.05}}
\put(4.75,3.75){\circle*{0.05}}

\end{picture}
\end{center}


\end{proof}

\end{document}
